




\section{Methods}

\subsection{Preliminaries: Exposure Bias in Autoregressive Video Diffusion Models}
An autoregressive video diffusion model is a hybrid generative framework that integrates autoregressive chain-rule decomposition with denoising diffusion for video generation.
Formally, given a sequence of $N$ video frames $x^{1:N} = (x^1, x^2, \dots, x^N)$, their joint distribution can be factorized using the chain rule: $p(x^{1:N}) = \prod_{i=1}^{N} p(x^i \mid x^{<i}).$ Each conditional distribution $p(x^i \mid x^{<i})$ is modeled through a diffusion process, where each frame is generated by progressively denoising Gaussian noise while conditioning on the previously generated frames. 
In practice, one may also generate a chunk of consecutive frames instead of a single frame at each step~\citep{yin2025slow, teng2025magi, huang2025self}. For clarity, we refer to each chunk simply as a frame in the following text.

Autoregressive video diffusion models are trained either (1) from scratch with frame-wise denoising loss or (2) by distilling a pretrained bidirectional model. The first approach is trained under the paradigm of Teacher Forcing (TF) or Diffusion Forcing (DF)~\citep{chen2024diffusion}. In TF, the conditional distribution for the $i$th frame at noise level $t_j$ is $p(x^i_{t_j} \mid x^{<i}_0)$, where all conditional history frames are the ground-truth clean frames from the training data. While in DF, the conditional distribution is $p(x^i_{t_j} \mid x^{<i}_{t_{\ge0}})$, where the history frames are the ground-truth frames corrupted with independent noise levels. Since training relies on ground-truth histories while inference relies on the model’s own predictions, a train–test gap known as exposure bias arises~\citep{schmidt2019generalization}. Mitigating the exposure bias is difficult because the denoising loss requires pairs of model predictions and the corresponding ground truth conditioned on them, which are unavailable. 

The second approach of distillation, however, provides a way to bypass the denoising loss and mitigate exposure bias. CausVid~\citep{yin2025slow} distills a pretrained bidirectional model into a few-step causal model. It adopts a Distribution Matching Distillation (DMD) loss~\citep{yin2024one} that minimizes the reverse KL divergence across randomly sampled timesteps $t$ between the smoothed data distribution $p_{\text{data}}(x_{t})$ and the student generator's output distribution $p_{\text{gen}}(x_{t})$. 
The gradient of the reverse KL can be approximated as the difference between two score functions:
\begin{multline}
\label{eq:dmd}
\nabla_{\theta} \mathcal{L}_{\text{DMD}} \triangleq \mathbb{E}_{t} \left( \nabla_{\theta} \text{KL} \left( p_{\text{gen}, t} \| p_{\text{data}, t} \right) \right) \\
 \approx - \mathbb{E}_{t} \left( \int \left( s_{\text{data}} \left( \Psi \left( G_{\theta}(\epsilon), t \right), t \right) - s_{\text{gen}} \left( \Psi \left( G_{\theta}(\epsilon), t \right), t \right) \right) \frac{d G_{\theta}(\epsilon)}{d \theta} \, d\epsilon \right), 
\end{multline}
where $\Psi$ represents the forward diffusion process, $\epsilon$ is random Gaussian noise, $G_{\theta}$ is the generator parameterized by $\theta$, and $s_{\text{data}}$ and $s_{\text{gen}}$ represent the score functions trained on the data and generator's output distribution, respectively. Since training with DMD loss does not require ground-truth image or video data~\citep{yin2024improved}, Self Forcing~\citep{huang2025self} mitigates the exposure bias by conditioning each frame on previously self-generated histories during training. However, although exposure bias is alleviated, severe error accumulation still occurs once generation extends beyond the trained temporal window.


\subsection{Autoregressive Video Generation via Rolling Diffusion Window}
\label{sec:inference}

\begin{figure}[t]
\centering
\vspace{-1em}
\includegraphics[width=\textwidth]{figure/method.pdf}
\caption{Illustration of the Rolling Forcing denoising process with $T=4$. Rolling Forcing jointly denoises a short window of consecutive frames that are assigned progressively higher noise levels and connected by bidirectional attention. The KV cache of recent frames is preserved as temporal context to maintain short-term consistency, while the KV cache of the initial frames is preserved as global context to ensure long-term consistency. During training, only a subset of denoising windows requires gradient computation, as highlighted by the red windows. These windows are mutually exclusive yet collectively cover all video frames.}
\label{fig:method}
\end{figure}

In Self Forcing (SF), videos are generated frame-by-frame in a strict causal manner. Consider a noise schedule $\{t_0=0,t_1, \dots, t_{T}=1000\}$ with total noise levels $T+1$. At each denoising step $t_j$ and frame index $i$, the model denoises an intermediate noisy frame \( x^i_{t_j} \) conditioned on previous clean frames \( x^{<i}_0 \) and then injects Gaussian noise with a lower noise level into the predicted denoised clean frame via the forward diffusion process $\Psi$. This produces a noisy frame \( x^i_{t_{j - 1}} \) which will be used as the input to the next denoising step. Formally, in SF, the denoising process is achieved by: $ x^i_{t_{j - 1}} = \Psi\big(G_\theta(x^i_{t_j}, t_{j},  x^{<i}_0), t_{j-1}\big)$, \mbox{and $ x^i_{t_{T}} \sim \mathcal{N}(0, I)$}. However, this formulation has no bidirectional attention between the current denoising frame $x^i$ and its history $x^{<i}$, where the strict causality forces every frame to inherit and compound the errors from its predecessors over time.



The proposed Rolling Forcing relaxes this constraint by extending the single-frame denoising window into a rolling window spanning multiple frames, as illustrated in \cref{fig:method}. Each denoising window contains consecutive frames with progressively higher noise levels in temporal order, akin to Rolling Diffusion~\citep{ruhe2024rolling}. The length of the denoising window $L_{\mathrm{win}}$ is set to the number of denoising time steps, \ie, $L_{\mathrm{win}} = T$. To ensure continuity, the next noise level of the $i$th frame is aligned with the current noise level of the $(i-1)$th frame, allowing the window to roll forward infinitely. At each roll, a clean frame is generated, and pure Gaussian noise is appended as the next frame to be synthesized. Formally, for the denoising window starting at the $i$th frame, the denoising distribution of Rolling Forcing can be defined by:
\begin{equation}
\label{eq:roll_window}
    p_\theta
    \Big(x^{i:i+T-1}_{t_{0:T-1}} \mid x^{i:i+T-1}_{t_{1:T}}, x^{<i}_0\Big) 
    = 
    \Psi\Big( G_\theta ( x^{i:i+T-1}_{t_{1:T}}, t_{1:T},  x^{<i}_0 ), t_{0:T-1}\Big),
\end{equation}
where $x^{i:i+T-1}_{t_{1:T}}$ denotes the noisy frames in the denoising window, and $x^{i:i+T-1}_{t_{0:T-1}}$ denotes the window output with each frame denoised to a lower noise level. The generator $G_\theta$ predicts clean frames conditioned on the input noisy frames, their noise levels $t_{1:T}$, and the clean history frames $x_0^{<i}$. $\Psi$ injects Gaussian noise $\epsilon_{t_{0:T-1}}$ at noise levels $t_{0:T-1}$ into the predicted clean frames, producing frames with reduced noise levels.

Since the length of the denoising window equals the number of denoising steps $T$, which is typically large (\ie, $\sim50$) in video diffusion models~\citep{wan2025wan}, the denoising window itself becomes prohibitively large. 
To manage such large windows, previous work either processes every frame independently on multiple GPUs~\citep{kim2024fifo}, or reduces $T$ to $\sim30$ using few-step samplers~\citep{xie2025progressive}.
In contrast, we adopt diffusion distillation~\citep{yin2024one, yin2024improved}, which reduces the number of denoising steps $T$ to just 5 while preserving generation quality, thereby making the denoising windows compact enough to fit on a single GPU while maintaining real-time latency.


\subsection{Temporal and Global History Context}
\label{sec:kvcache}
% \begin{figure}[t!]

\begin{wrapfigure}{r}{0.48\textwidth}
\vspace{-2em}
\begin{minipage}[t]{0.48\textwidth}
  \begin{algorithm}[H]
    \caption{Rolling Forcing Training}
    \small
    \begin{algorithmic}[1]
      \Require Denoise timesteps $\{t_0, t_1, \dots, t_T\}$
      \Require Number of video frames $N$
      \Require AR diffusion model $G_\theta$ (returns KV embeddings via $G_\theta^{\mathrm{KV}}$)
      \Loop
        \State Initialize model output $\mathbf{X}_{\theta} \gets []$
        \State Initialize KV cache $\mathbf{KV} \gets []$
        \State Initialize $x^{1:T-1}_{t_{1:T-1}}$ with $G_\theta$
        \State Sample $j \sim \text{Uniform}\{0, 1, \ldots, T-1\}$
        \For{$i = 1, \dots, N$}
          \State Sample $x_{t_T}^{i+T-1} \sim \mathcal{N}(0, I)$
          \State Set $x^{i:i+T-1}_{t_{1:T}} \gets x^{i:i+T-2}_{t_{1:T-1}} \,\|\, x_{t_T}^{i+T-1}  $
          \State Select and apply RoPE to $\mathbf{KV}$ (\cref{sec:kvcache})
          \If{$i \equiv j \pmod{T}$}
            \State Enable gradient computation
            \State $\hat{x}^{i:i+T-1}_0 \gets G_\theta(x^{i:i+T-1}_{t_{1:T}}, t_{1:T},  \mathbf{KV})$
            \State $\mathbf{X}_{\theta}{\texttt{.append}}( \hat{x}^{i:i+T-1}_0)$
            \State Disable gradient computation
          \Else 
            \State $\hat{x}^{i:i+T-1}_0 \gets G_\theta(x^{i:i+T-1}_{t_{1:T}}, t_{1:T},  \mathbf{KV})$
          \EndIf
          \State $\mathbf{KV}{\texttt{.append}}( G_\theta^\text{KV}(\hat{x}^i_{0}, t_0, \mathbf{KV}))$
          \State $x^{i+1:i+T-1}_{t_{1:T-1}} \gets \Psi(\hat{x}^{i+1:i+T-1}_0, t_{1:T-1}) $
        \EndFor
        \State Update $\theta$ via DMD loss (\cref{eq:dmd})
      \EndLoop
    \end{algorithmic}
    \label{alg:training}
  \end{algorithm}
\end{minipage}
\vspace{-2em}
\end{wrapfigure}

% \end{figure}






As the clean history frames $x^{<i}_0$ accumulate during generation, handling them directly becomes computationally expensive. To address this, following~\citet{huang2025self}, we cache the key and value states of the history frames, thereby avoiding redundant recomputation when generating new frames, as illustrated in \cref{fig:method}. Note that although the attention within the denoising window is bidirectional, the attention between the frames in the denoising window and the KV cache of history frames remains causal. 
While KV caching reduces computation, the computational complexity still grows quadratically with the cache size as frames accumulate, and the cache may become large enough to cause out-of-memory errors. Given a denoising window starting at the $i$th frame $x^{i:i+T-1}_{t_{1:T}}$, we address this issue by retaining only the KV cache of the most recent $L_{\mathrm{tem}}$ history frames $x^{i-L_{\mathrm{tem}}:i-1}_0$ as temporal context to preserve short-term temporal consistency. However, relying solely on short-term history causes a gradual drift of long-range properties of the generated video (like exposure, color tone, white balance, etc.) as generation proceeds.



To maintain long-term global consistency, we cache the KV states of the initial $L_{\mathrm{glo}}$ generated frames $x_0^{1:L_{\mathrm{glo}}}$ as global context, analogous to attention sink tokens in streaming language models~\citep{xiao2023efficient}. The cache sizes $L_{\mathrm{tem}}$ and $L_{\mathrm{glo}}$ are chosen such that the total attention window size matches that of the bidirectional teacher model, \ie, $L_{\mathrm{tem}} + L_{\mathrm{glo}} + L_{\mathrm{win}} = L_{\mathrm{bidirectional}}$.
However, directly caching the initial frames leads to spilling problems. Modern video diffusion DiTs~\citep{peebles2023scalable} typically use RoPE~\citep{su2024roformer} for relative positional encoding. As the indices of the denoising frames $i:i+T-1$ increase, their relative distance to the initial cached frames grows, eventually exceeding the trained range of RoPE and producing unnatural artifacts.
To resolve this, we cache the key states of the global context frames $x_0^{1:L_{\mathrm{glo}}}$ before applying the RoPE transformation. During generation, we dynamically apply RoPE to these cached key states at the effective indices $  i - L_{\mathrm{tem}} - L_{\mathrm{glo}} :i - L_{\mathrm{tem}}-1$, treating them as being positioned immediately before the temporal context frames $x^{i-L_{\mathrm{tem}}:i-1}_0$. This adjustment preserves a fixed relative position \wrt the denoising frames, preventing excessive offsets.


\subsection{Rolling Forcing Post-Training}
\label{sec:training}




Rolling Forcing distills a pretrained bidirectional video diffusion model~\citep{wan2025wan} to a few-step causal autoregressive generator using the DMD loss (\cref{eq:dmd}). As DMD matches the holistic distribution of the entire video sequence to the data distribution $D(p_{\text{data}}(x^{1:N}) \| p_\theta(x^{1:N}))$, the calculation of the DMD loss requires a predicted clean video $\hat{x}_0^{1:N}$ during training. In SF, the predicted clean video is generated by: 
\begin{equation}
\label{eq:sf}
 \hat{x}_0^{1:N} = \Big\{ 
 \hat{x}^i_0= G_\theta(x^{i}_{t_j}, t_j, x^{<i}_0) 
 \mid i=1,2,\dots,N
 \Big\},
\end{equation}
where $j \sim \text{Uniform}\{0,1,\dots,T-1\}$ indicates each frame's noise level $t_j$ before denoising.
For Rolling Forcing, as the denoising window consists of multiple frames at different noise levels, we select the $j$th frame in each window and combine the selected frames as the predicted clean video:
\begin{equation}
\label{eq:rf_init}
    \hat{x}_0^{1:N} = \Big\{ 
    \hat{x}^i_0 = (\hat{x}^{i:i+T-1}_0)^j = 
    G_\theta(x^{i:i+T-1}_{t_{1:T}}, t_{1:T},  x^{<i}_0)^j 
    \mid i=1,2,\dots,N
    \Big\}, 
\end{equation}
where $j \sim \text{Uniform}\{0,1,\dots,T-1\}$ represents both the frame's index within the denoising window and the frame's noise level $t_j$. However, \cref{eq:rf_init} incurs $T$ times higher computational complexity than \cref{eq:sf}, because the query size is $T$ times larger. Given that DMD loss is already computationally expensive, this additional cost can easily lead to out-of-memory error even on GPUs with 80G of memory. 


To address this issue, instead of backpropagating through every window (which requires gradients for each forward pass), we sample a subset of non-overlapping windows to construct the predicted clean video, as illustrated in \cref{fig:method}. Gradient computation is performed only on these selected windows, significantly reducing memory usage while retaining effective supervision. Formally, the predicted clean video is given by:
\begin{equation}
\label{eq:rf_final}
    \hat{x}_0^{1:N} = \Big\{ \hat{x}^{i:i+T-1}_0 = 
    G_\theta(x^{i:i+T-1}_{t_{1:T}}, t_{1:T},  x^{<i}_0) \mid 
    i \equiv j \pmod{T},\ 1 \leq i \leq N
    \Big\}, 
\end{equation}
where $j \sim \text{Uniform}\{0,1,\dots,T-1\}$. In each iteration, we reduce the number of forward passes requiring gradient computation from $N$ in \cref{eq:rf_init} to $\left\lceil N/T \right\rceil$\footnote{We omit the denoising windows at the start of the video for clarity in \cref{eq:rf_final,eq:roll_window}, where the window has fewer than $T$ frames. Gradient computation is still required if the window index $i$ satisfies $i \equiv j \pmod{T}$.}. The Rolling Forcing training is illustrated in \cref{alg:training}. Similar to SF, the input noisy frames $x^{i:i+T-1}_{t_{1:T}}$ during training are generated by the model rather than taken from ground truth, thus mitigating the exposure bias. However, unlike \cref{eq:sf} or \cref{eq:rf_init}, where every frame in the predicted clean video $\hat{x}_0^{1:N}$ is denoised from the same noise level $t_j$, the frames in \cref{eq:rf_final} are denoised from varying noise levels $t_{1:T}$. Consequently, frames denoised from different noise levels have different quality and clarity, leading to unnatural video $\hat{x}_0^{1:N}$ and camera movement in DMD training. To address this issue, we adopt a mixed training strategy that alternates between SF training (\cref{eq:sf}) and Rolling Forcing training (\cref{eq:rf_final}) with equal probability. The SF objective serves as a regularizer, encouraging the model to produce videos with natural camera movement. The inference adopts the Rolling Forcing paradigm alone as elaborated in \cref{alg:inference}.



